% !TeX root = ../../金融经济学_王江_讲义及习题.tex
% 教材练习：印刷第 82—85 页（PDF 第 97—100 页）；答案参考 PDF 第 20—25 页。
% 5.1 教材原题除息价为 uS、dS；答案参考题干改写为 (1+rF+u)S、(1+rF-d)S 却仍沿用 uS、dS 的解答。恢复教材原题。
% 5.3 原题“较高”补明弱不等式及严格条件；有限状态下两支付可同时为零。
% 5.4 原参考答案用当期内在价值差界定未来蝴蝶支付价格，一般不成立；以到期支付上界 delta 贴现。
% 5.8(c) 原题严格较高改为不低于；新增可执行策略只保证弱不等式。
% 5.9 原参考仅解 ud<1 的子情形；现对四条路径取最大值，覆盖题给全部参数。
\newcounter{wjFiveExercise}
\renewcommand{\thewjFiveExercise}{\thechapter.\arabic{wjFiveExercise}}
\newcommand{\WJFiveProblem}[1]{\refstepcounter{wjFiveExercise}\label{prob:wj-five-\arabic{wjFiveExercise}}\problemno{\thewjFiveExercise#1}}

\begin{problem}
\WJFiveProblem{} 1 期有两个可能状态 $a$ 和 $b$，发生概率分别为 $\pi$ 和 $1-\pi$。存在一只收益率为 $r_{F}$ 的无风险证券。还存在一只风险证券，它的现时价格为 $S$，在 1 期支付红利 $\delta S$。它在 1 期的除权价（即红利后的价格）在状态 $a$ 下是 $uS$、在状态 $b$ 下是 $dS$。
\begin{enumerate}[label=(\alph*)]
    \item 为使市场中不存在套利机会，$\pi$，$r_{F}$，$u$，$d$ 和 $\delta$ 应该满足什么条件？
    \item 1 期到期、执行价为 $K$ 的欧式买权的 0 期价格是多少？
    \item 美式买权的价格是多少？比较它与相应的欧式买权价格。什么时候提前执行是最优的？
    \item 将买权换成卖权，回答 (b) 和 (c)。
\end{enumerate}
\end{problem}
\begin{solution}

令 $R=1+r_F$。按 $u>d\ge0$、$S>0$、$\delta\ge0$ 记号讨论；期权在 1 期股利支付后按除息价结算。股票持有人收到的总支付分别为 $(u+\delta)S$、$(d+\delta)S$。
\begin{enumerate}[label=(\alph*)]
 \item 取两个状态均有正概率，即 $0<\pi<1$，并要求 $R>0$。无套利的充要条件为
 \[
 d+\delta<R<u+\delta.
 \]

 若 $d+\delta\ge R$，买入一股并借入 $S$，到期两状态净支付均非负，至少一项严格为正；若 $u+\delta\le R$，卖出一股并将所得存入债券账户，也得到套利。反过来，下列严格正状态价格同时为股票和债券定价，保证无套利。

 \item 状态价格满足
 \[
 \phi_a+\phi_b=\frac1R,\qquad
 (u+\delta)S\phi_a+(d+\delta)S\phi_b=S.
 \]

 解得
 \[
 \phi_a=\frac{R-d-\delta}{R(u-d)},\qquad
 \phi_b=\frac{u+\delta-R}{R(u-d)}.
 \]

 欧式看涨在除息后结算，故
 \[
 c=\phi_a[uS-K]_++\phi_b[dS-K]_+.
 \]

 具体为
 \[
 c=\begin{cases}
 0,&K\ge uS,\\
 \phi_a(uS-K),&dS\le K<uS,\\
 \phi_a(uS-K)+\phi_b(dS-K),&K<dS.
 \end{cases}
 \]

 股票在到期时还支付确定股利 $\delta S$，所以欧式平价可核验为
 \[
 c+\frac K R+\frac{\delta S}R=p+S.
 \]

 \item 美式看涨在 0 期立即执行并卖出股票可得 $[S-K]_+$，等待则价值为 $c$，因此
 \[
 C=\max\{[S-K]_+,c\}\ge c.
 \]

 当 $[S-K]_+>c$ 时立即执行严格更优；等号时两种选择价值相同。固定 $S,u,d,R$，在无套利参数范围内增加 $\delta$ 会降低 $\phi_a$、提高 $\phi_b$。因看涨的上涨支付不低于下跌支付，继续持有价值 $c$ 非增，提前执行动机相应增强。

 \item 欧式看跌为
 \[
 p=\phi_a[K-uS]_++\phi_b[K-dS]_+
 =\begin{cases}
 0,&K\le dS,\\
 \phi_b(K-dS),&dS<K\le uS,\\
 \phi_a(K-uS)+\phi_b(K-dS),&K>uS.
 \end{cases}
 \]

 美式看跌为 $P=\max\{[K-S]_+,p\}$，立即执行严格更优的条件为 $[K-S]_+>p$。看跌的下跌支付不低于上涨支付；在上述固定参数比较中，增加 $\delta$ 使继续持有价值 $p$ 非减，从而减弱提前执行动机。
\end{enumerate}
\end{solution}

\begin{problem}
\WJFiveProblem{} 在 1 期有三个可能状态：$\{\omega_{1},\omega_{2},\omega_{3}\}$，在经济中有三只交易证券，它们的支付向量如下：
\[
    X_{a}=\begin{bmatrix}1\\1\\0\end{bmatrix},\quad
    X_{b}=\begin{bmatrix}0\\1\\1\end{bmatrix},\quad
    X_{c}=\begin{bmatrix}1/2\\1\\1/2\end{bmatrix}.
\]

这些证券的 0 期均衡价格是 $S_{a}=S_{b}=S_{c}=1$。令 $\phi=\begin{bmatrix}\phi_{1} \\ \phi_{2} \\ \phi_{3}\end{bmatrix}$ 为状态价格向量。考虑一份标的资产证券为 $a$、执行价格为 $1/2$ 的欧式买权。它的支付是什么？你能确定它的价格吗？
\end{problem}
\begin{solution}

它的支付是 $\begin{bmatrix}1/2\\1/2\\0\end{bmatrix}$，因为它的支付恰为证券 $a$ 的一半，因而它的价格也是一半即 $\tfrac12$。
\end{solution}

\begin{problem}
\WJFiveProblem{} 蝴蝶头寸是看涨期权的组合：（1）买入 1 份执行价格为 $K-\delta$ 的看涨期权，（2）卖出 2 份执行价格为 $K$ 的看涨期权，（3）买入 1 份执行价格为 $K+\delta$ 的看涨期权。因此，蝴蝶头寸可由两个参数 $K$ 和 $\delta$ 刻划。比较两个蝴蝶头寸，$K$ 相同但 $\delta$ 不同。证明较宽的蝴蝶头寸价值不低，并说明严格较高的条件。
\end{problem}
\begin{solution}

蝴蝶支付可简写为
\[
 h_\delta(x)=[\delta-|x-K|]_+.
\]

若 $\delta_1>\delta_0>0$，逐状态均有 $h_{\delta_1}(x)\ge h_{\delta_0}(x)$。由\ifdefined\WJHandoutKnowledge 定理~\ref{thm:wj-four-5}\else 价格的单调性\fi（逐状态支付不低则价格不低），较宽组合的价格不低。

当存在正概率状态满足 $|\tilde X-K|<\delta_1$ 时，较宽组合在该状态的支付严格更高，由\ifdefined\WJHandoutKnowledge 定理~\ref{thm:wj-four-4}\else 非负支付的正价格性质\fi（非负、非零支付具有严格正价格）得到严格价格比较。若全部状态均落在较宽组合的支付区间之外或端点，两组合支付都为零，价格相等。
\end{solution}

\begin{problem}
\WJFiveProblem{} 对于参数为 $K$ 和 $\delta$ 的蝴蝶头寸，给出它的价格的一个上界和下界。
\end{problem}
\begin{solution}

蝴蝶支付满足 $0\le h_\delta(\tilde X)\le\delta$。若单位债券到期支付 $1$、价格为 $B=1/R$，则无套利给出
\[
 0\le V(h_\delta)\le\frac\delta R.
\]

下界由支付非负得到；上界由到期确定支付 $\delta$ 的债券组合得到。还可利用
\[
 h_\delta(\tilde X)\le[\tilde X-(K-\delta)]_+-[\tilde X-K]_+
\]
给出可交易看涨价差的上界：
\[
 V(h_\delta)\le c(S,K-\delta)-c(S,K).
\]
\end{solution}

\begin{problem}
\WJFiveProblem{} 股票价格具有如下形式的二叉树结构：
\begin{center}
\WJEventTree[left={$10$},stem]{12,8}
\end{center}
两个可能结果的发生概率相等，无风险利率是 $10\%$。
\begin{enumerate}[label=(\alph*)]
    \item 为一个以此股票为标的资产、执行价格为 9 的看跌期权定价。
    \item 假设你得到一个新的利好信息：股票价格上升到 12 的概率大于 $1/2$，而股票价格仍然是 10。你愿意支付的看跌期权的价格是上涨还是下跌？解释原因。
\end{enumerate}
\end{problem}
\begin{solution}

\begin{enumerate}[label=(\alph*)]
 \item $R=11/10$、$u=6/5$、$d=4/5$，故风险中性上涨概率为 $q=(R-d)/(u-d)=3/4$。状态价格为
 \[
 \phi_u=\frac qR=\frac{15}{22},\qquad
 \phi_d=\frac{1-q}R=\frac5{22}.
 \]

 看跌到期支付为 $(0,1)^{\mathsf T}$，所以 $p_0=5/22$。
 \item 给定股票和债券价格、未来两种支付，复制组合与状态价格已确定。真实上涨概率的变化不改变这组复制条件，故看跌的无套利价格仍为 $5/22$。
\end{enumerate}
\end{solution}

\begin{problem}
\WJFiveProblem{} 二叉树模型可以拓展到多于两个时期的情形。考虑拓展到三个时期 0，1，2 的情况。股票价格服从如下的二叉树过程：
\begin{center}
\begin{tikzpicture}[WJ binomial compact]
\node[WJ binomial node] (A) at (0,0) {$10$};
\node[WJ binomial node] (B) at (1,1.25) {$12$};
\node[WJ binomial node] (C) at (1,-1.25) {$8$};
\node[WJ binomial node] (D) at (2,2) {$14.4$};
\node[WJ binomial node] (E) at (2,.6667) {$9.6$};
\node[WJ binomial node] (F) at (2,-.6667) {$9.6$};
\node[WJ binomial node] (G) at (2,-2) {$6.4$};
\foreach \a/\b in {A/B,A/C,B/D,B/E,C/F,C/G}
  \draw[WJ binomial edge] (\a.east) -- (\b.west);
\foreach \x/\t in {0/0,1/1,2/2}
  \node[WJ binomial time] at (\x,-2.65) {$t=\t$};
\end{tikzpicture}
\end{center}
每一期间（即，从 0 期到 1 期或从 1 期到 2 期）的利率都是 $10\%$。考虑一份在 2 期到期、执行价格为 9 的看涨期权。
\begin{enumerate}[label=(\alph*)]
    \item 描述这份看涨期权在 2 期的支付。
    \item 求解看涨期权在 1 期、股票价格分别为 12 或 8 时的价格。
    \item 求解看涨期权在 0 期的价格。
    \item 如果是美式看涨期权，它在 0 期的价格是多少？
\end{enumerate}
\end{problem}
\begin{solution}

沿用练习~\ref{prob:wj-five-5} 的逐期状态价格 $\phi_u=15/22$、$\phi_d=5/22$，每个节点使用相同的复制关系。
\begin{enumerate}[label=(\alph*)]
 \item 按 $uu,ud,du,dd$ 四条路径排列，到期支付为 $(27/5,3/5,3/5,0)^{\mathsf T}$。
 \item 股价为 $12$ 与 $8$ 时的期权价格分别为
 \[
 c_{1,u}=\frac{15}{22}\frac{27}{5}+\frac5{22}\frac35=\frac{42}{11}\approx3.8182,
 \qquad
 c_{1,d}=\frac{15}{22}\frac35=\frac9{22}\approx0.4091.
 \]
 \item 继续倒推，
 \[
 c_0=\frac{15}{22}\frac{42}{11}+\frac5{22}\frac9{22}
 =\frac{1305}{484}\approx2.6963.
 \]
 \item 股票无股利且利率为正，各节点提前执行看涨均不优于持有。美式与欧式价格相同，约为 $2.70$。
\end{enumerate}
\end{solution}

\begin{problem}
\WJFiveProblem{*} 在练习~\ref{prob:wj-five-6} 描述的经济中，考虑一份以股票为标的资产、在 2 期到期、执行价格为 14 的看跌期权。
\begin{enumerate}[label=(\alph*)]
    \item 如果它是欧式期权，求解它的 0 期价格。
    \item 如果它是美式期权，求解它的 0 期价格。
\end{enumerate}
\end{problem}
\begin{solution}

\begin{enumerate}[label=(\alph*)]
 \item 到期支付为 $(0,22/5,22/5,38/5)^{\mathsf T}$。在 1 期两节点，
 \[
 p_{1,u}=\frac5{22}\frac{22}{5}=1,\qquad
 p_{1,d}=\frac{15}{22}\frac{22}{5}+\frac5{22}\frac{38}{5}
 =\frac{52}{11}.
 \]

 因而欧式价格为
 \[
 p_0=\frac{15}{22}\times1+\frac5{22}\frac{52}{11}
 =\frac{425}{242}\approx1.7562.
 \]
 \item 美式期权在 1 期比较立即执行与继续持有：
 \[
 P_{1,u}=\max\{14-12,1\}=2,\qquad
 P_{1,d}=\max\{14-8,52/11\}=6.
 \]

 两节点均应立即执行。0 期继续持有价值为 $(15/22)2+(5/22)6=30/11$，立即执行价值为 $14-10=4$，故 $P_0=4$。
\end{enumerate}
\end{solution}

\begin{problem}
\WJFiveProblem{*} 这里的经济与练习~\ref{prob:wj-five-6} 描述的经济一样，只是这里的参数更一般性：
\begin{center}
\begin{tikzpicture}[WJ binomial compact]
\node[WJ binomial node] (A) at (0,0) {$S$};
\node[WJ binomial node] (B) at (1,1.25) {$uS$};
\node[WJ binomial node] (C) at (1,-1.25) {$dS$};
\node[WJ binomial node] (D) at (2,2) {$u^{2}S$};
\node[WJ binomial node] (E) at (2,.6667) {$udS$};
\node[WJ binomial node] (F) at (2,-.6667) {$duS$};
\node[WJ binomial node] (G) at (2,-2) {$d^{2}S$};
\foreach \a/\b in {A/B,A/C,B/D,B/E,C/F,C/G}
  \draw[WJ binomial edge] (\a.east) -- (\b.west);
\foreach \x/\t in {0/0,1/1,2/2}
  \node[WJ binomial time] at (\x,-2.65) {$t=\t$};
\end{tikzpicture}
\end{center}
每一期间的利率都是 $r_{F}$ 且 $0<d<1+r_{F}<u$。
\begin{enumerate}[label=(\alph*)]
    \item 求解在每一时期考虑下一期的两种可能情形（即股价上升或下降）的风险中性概率。
    \item 用 $S$，$u$，$d$ 和 $r_{F}$ 表示欧式看跌期权在 0 期的价格。
    \item 证明在 2 期到期的美式看涨期权的价格比具有相同执行价格、但在 1 期到期的美式看涨期权的价格不低。
\end{enumerate}
\end{problem}
\begin{solution}

令 $R=1+r_F$。由 $0<d<R<u$，每个节点的风险中性概率均为
\[
 q=\frac{R-d}{u-d},\qquad1-q=\frac{u-R}{u-d}.
\]
\begin{enumerate}[label=(\alph*)]
 \item 下一期条件状态价格为 $\phi_u=q/R$、$\phi_d=(1-q)/R$，各节点相同。
 \item 四条路径的支付按其风险中性权重贴现两期：
 \[
 p_0=\frac1{R^2}\left\{
 q^2[K-u^2S]_++2q(1-q)[K-udS]_++(1-q)^2[K-d^2S]_+
 \right\}.
 \]
 \item 2 期到期的美式期权允许在 0、1、2 期执行，1 期到期的只允许在 0、1 期执行。后者的每一种执行策略都可用于前者，所得支付相同，故前者的价值不低。

 新增执行机会可能带来严格增益。例如 $S=10,u=1.2,d=0.8,R=1.1,K=13$ 时，1 期到期的看涨在两状态均为零，而 2 期到期的看涨在 $uu$ 路径有正支付，因此后者价格严格更高。若把 $K$ 改为 $20$，所有执行时点都为虚值，两种期限的价格均为零。
\end{enumerate}
\end{solution}

\begin{problem}
\WJFiveProblem{*} 考虑练习~\ref{prob:wj-five-8} 描述的经济。假设一只证券在 2 期的支付是股票价格从 0 期到 2 期所达到的最大值。（这是一个关于回望期权的例子，它的支付依赖于标的证券的价格路径）。
\begin{enumerate}[label=(\alph*)]
    \item 写出它在 2 期的支付。
    \item 导出它的 0 期价格。
    \item 证明用在 2 期支付一美元的无风险债券的价格正规化以后，这只证券的价格过程在等价鞅测度下是鞅。
\end{enumerate}
\end{problem}
\begin{solution}

令 $R=1+r_F$、$q=(R-d)/(u-d)$。记 $h_{ij}$ 为沿路径 $ij$ 到达 2 期时的最大历史股价。
\begin{enumerate}[label=(\alph*)]
 \item 按 $uu,ud,du,dd$ 排列，完整支付为
 \[
 \begin{aligned}
 h_{uu}&=S\max\{1,u,u^2\},&h_{ud}&=S\max\{1,u,ud\},\\
 h_{du}&=S\max\{1,d,du\},&h_{dd}&=S\max\{1,d,d^2\}.
 \end{aligned}
 \]

 若 $u>1>d$ 且 $ud<1$，则得到 $(u^2S,uS,S,S)^{\mathsf T}$。中间两条路径的终点股价相同，但历史最高价可以不同。

 \item 1 期的两个节点价值与 0 期价值分别为
 \[
 V_{1,u}=\frac{qh_{uu}+(1-q)h_{ud}}R,\qquad
 V_{1,d}=\frac{qh_{du}+(1-q)h_{dd}}R,
 \]
 \[
 V_0=\frac1{R^2}\left\{q^2h_{uu}+q(1-q)(h_{ud}+h_{du})+(1-q)^2h_{dd}\right\}.
 \]

 四条路径分别赋予风险中性权重；相同终点股价的路径按各自的历史最高价计价。

 \item 以 2 期支付一单位的债券为计价单位，其价格为
 \[
 B_0=R^{-2},\qquad B_1=R^{-1},\qquad B_2=1.
 \]

 记 $\widehat V_t=V_t/B_t$。在 1 期上涨节点，
 \[
 \widehat V_{1,u}=RV_{1,u}=qh_{uu}+(1-q)h_{ud}
 =\left.\mathrm E^Q[\widehat V_2\mid\mathcal F_1]\right|_u.
 \]

 在下跌节点同样有 $\widehat V_{1,d}=qh_{du}+(1-q)h_{dd}$。0 期则有
 \[
 \mathrm E^Q[\widehat V_1]
 =qRV_{1,u}+(1-q)RV_{1,d}=R^2V_0=\widehat V_0.
 \]

 两个日期的条件期望关系均成立，所以债券计价的价格过程在等价鞅测度 $Q$ 下是鞅。
\end{enumerate}
\end{solution}

\begin{problem}
\begin{samepage}
\WJFiveProblem{} 经济在 1 期有 3 个状态。状态指数证券在三个状态下的支付分别为 1，2，3：
\begin{center}
\WJEventTree[stem]{1,2,3}
\end{center}
并有两只以指数证券为标的资产、执行价格分别为 1 和 2 的看涨期权。用指数证券和上述期权复制如下支付：
\begin{center}
\WJEventTree[stem]{1,1,1}
\qquad
\WJEventTree[stem]{3,2,1}
\qquad
\WJEventTree[stem]{1,0,0}
\end{center}
\end{samepage}
\end{problem}
\begin{solution}

记状态指数证券的支付向量为 $X_{0}=\begin{bmatrix}1\\2\\3\end{bmatrix}$，执行价格为 1 的指数期权的支付向量为 $X_{1}=\begin{bmatrix}0\\1\\2\end{bmatrix}$，执行价格为 2 的指数期权的支付向量为 $X_{2}=\begin{bmatrix}0\\0\\1\end{bmatrix}$ 那么我们可得：
\[
    X_{0}-X_{1}=\begin{bmatrix}1\\1\\1\end{bmatrix},\quad
    3X_{0}-4X_{1}=\begin{bmatrix}3\\2\\1\end{bmatrix},\quad
    X_{0}+X_{2}-2X_{1}=\begin{bmatrix}1\\0\\0\end{bmatrix}.
\]
\end{solution}

